
\begin{figure}[ht!]
    \centering
    \footnotesize
    \begin{tikzpicture}
        \begin{axis}[
            ybar,
            clip=false,
            symbolic x coords={
                Low\\On,
                Low\\Off,
                High\\On,
                High\\Off
            },
            xtick=data,
            ymin=0,
            ymax=240,
            ylabel={Latência Média (ms)},
            bar width=20pt,
            enlarge x limits=0.15,
            xticklabel style={align=center},
            width=12cm,
            height=7cm,
            legend style={at={(0.5,1.05)}, anchor=south, legend columns=-1}
        ]
        % SMMS (green)
        \addplot+[
            fill=mygreen,
            draw=black,
            nodes near coords,
            every node near coord/.append style={font=\footnotesize, yshift=7pt, text=mygreen},
            error bars/.cd,
            y dir=both,
            y explicit,
            error bar style={black}
        ] coordinates {
            (Low\\On, 20.0) +- (0.0,0.0)
            (Low\\Off, 20.0) +- (0.0,0.0)
            (High\\On, 60.0) +- (0.0,0.0)
            (High\\Off, 26.0) +- (0.0,0.0)
        };

        % Thea (red)
        \addplot+[
            fill=myred,
            draw=black,
            nodes near coords,
            every node near coord/.append style={font=\footnotesize, yshift=7pt, text=myred},
            error bars/.cd,
            y dir=both,
            y explicit,
            error bar style={black}
        ] coordinates {
            (Low\\On, 170.97) +- (9.64,9.64)
            (Low\\Off, 28.7) +- (3.01,3.01)
            (High\\On, 205.98) +- (0.28,0.28)
            (High\\Off, 28.4) +- (1.26,1.26)
        };
        
        % EDF (blue)
        \addplot+[
            fill=myblue,
            draw=black,
            nodes near coords,
            every node near coord/.append style={font=\footnotesize, yshift=7pt, text=myblue},
            error bars/.cd,
            y dir=both,
            y explicit,
            error bar style={black}
        ] coordinates {
            (Low\\On, 81.04) +- (4.85,4.85)
            (Low\\Off, 29.01) +- (1.4,1.4)
            (High\\On, 136.69) +- (3.46,3.46)
            (High\\Off, 29.61) +- (2.81,2.81)
        };

        \legend{SMMS, Thea, EDF}

        % Custom label on the left
        \draw
          (axis description cs:-0.165,-0.19)
          node[anchor=south west, align=left] {\footnotesize
            \textbf{Workload:}\\
            \textbf{Nuvem:}
          };
        \end{axis}
    \end{tikzpicture}
    \caption{Desempenho médio de latência em função da carga de trabalho e da configuração da nuvem.}
    \label{fig:results2}
\end{figure}

A Figura \ref{fig:results2} ilustra a latência média das requisições em cada cenário para os três algoritmos. Nota-se que o SMMS alcançou as menores latências médias em todos os cenários. Sob carga baixa, a latência média do SMMS ficou em torno de 20 ms (tanto com nuvem ligada quanto desligada), indicando que o SMMS conseguiu atender as requisições majoritariamente na borda local, mantendo o tempo de resposta baixo e praticamente igual ao mínimo (latência intrínseca da borda). O EDF apresentou latências maiores (por exemplo, \textasciitilde81 ms em Low On), e o Thea teve as latências médias mais elevadas quando a nuvem estava disponível: aproximadamente 171 ms em Low On e 206 ms em High On.

Esses valores altos para o Thea refletem sua estratégia de offloading, resultando em muitas requisições sendo enviadas para a nuvem e aumentando substancialmente o tempo médio de resposta. Em contraste, o SMMS equilibra melhor esses fatores e evita usar a nuvem desnecessariamente, o que lhe permite manter latências próximas das mínimas possíveis. Nos cenários sem nuvem (Low Off e High Off), todos os algoritmos apresentam latências médias similares e baixas (em torno de 20–30 ms), já que todas as requisições são servidas nas bordas. Um ponto interessante é que, em carga alta com nuvem ativada (High On), o SMMS obteve latência média em torno de 60 ms, consideravelmente menor que a do Thea (206 ms) e do EDF (137 ms). Isso indica que, mesmo utilizando a nuvem quando necessário em alta carga, o SMMS consegue limitar o impacto no tempo de resposta médio, possivelmente enviando apenas o estritamente necessário para a nuvem e mantendo a maior parte do processamento próximo dos usuários.

